\documentclass[a4paper,10pt]{jsarticle}

% --- パッケージの読み込み ---
\usepackage[margin=25mm]{geometry} % 上下左右の余白を25mmに設定
\usepackage{amsmath,amssymb}       % 高度な数式環境・記号
\usepackage[dvipdfmx]{graphicx}    % 画像挿入用（LuaLaTeXの場合は [dvipdfmx] を削除）
\usepackage{url}                   % URL表示用

% --- ヘッダー・フッターの設定 ---
\pagestyle{plain}                  % ページ下部中央のみにページ番号を表示

\begin{document}

% -----------------------------------------------------------------
\section{目的}
ガイガー・ミュラー計数管（以下, GM計数管と記す）を用いて、放射性元素の崩壊の法則と物質による放射線の吸収について調べ、崩壊の際に発生する線を用いて、金属箔による電子線の吸収係数を求める。

% -----------------------------------------------------------------
\section{原理}

\subsection{放射性元素の崩壊}
原子核は陽子と中性子からなっている。しかし、原子核の質量は、陽子と中性子をバラバラにした質量の和より小さく、この質量の差のことを質量欠損という。原子核の中には不安定なものが存在し、不安定な原子核は崩壊してより安定な原子核になる。ここで、線や線や線を放出して、より安定な別の原子核になる現象を放射性元素の崩壊とよぶ。

\subsection{放射線崩壊の計数値の分布}
放射性原子核の崩壊は確率論的な現象であり、個々の原子核がいつ崩壊するかを正確に予測することはできない。時刻 t における放射性核種の数を N とすると、微小時間 dt の間に崩壊する核種数 dN は N に比例するため、以下の関係式が成り立つ。
\begin{equation}
  \frac{dN}{dt} = -\lambda N
\end{equation}
ここで  は崩壊定数（>0）である。時刻 t=0 での核種数を N0 とすると、時刻 t において崩壊せずに残っている平均核種数 N(t) は次のように表される。
\begin{equation}
  N(t) = N_0 e^{-\lambda t}
\end{equation}

半減期を t1/2 とすると、e-t1/2=1/2 より t1/2=20.693 となる。t1 の近似が成り立つ短い測定時間 t においては、e-t1-t とテーラー展開できるため、時間 t 内に崩壊する期待値（平均値）N は次のように近似される。
\begin{equation}
  \bar{N} = N_0 \lambda t
\end{equation}

一定時間内に観測される崩壊数 N の出現確率 P(N) は、平均値を N とするポアソン分布に従う。
\begin{equation}
  P(N) = \frac{\bar{N}^N}{N!} e^{-\bar{N}} \quad \left( \sum_{N=0}^{\infty} P(N) = 1 \right)
\end{equation}
ポアソン分布における標準偏差  は、平均値の平方根に等しく =N となる。

さらに、平均計数値 N が十分に大きい場合には、ポアソン分布は中心極限定理により次のようなガウス分布（正規分布） G(N) に近似される。
\begin{equation}
  G(N) = \frac{1}{\sqrt{2\pi \bar{N}}} \exp \left\{ -\frac{(N - \bar{N})^2}{2\bar{N}} \right\}
\end{equation}
本実験では、ゲート時間内における計数値の出現頻度を記録し、その分布がポアソン分布（またはガウス分布）に適合するか検証を行う。

\subsection{線の吸収}
線（高速の電子）が物質中を通過する際、物質内の原子や電子と衝突して励起や電離を引き起こし、運動エネルギーを失いながら徐々に減衰する。厚さ x の物質層に入射する前の電子数を N0、層を通過した後の電子数を N(x) とすると、線の強さは物質の厚さに対して指数関数的に減衰し、次の関係式が成り立つ。
\begin{equation}
  N(x) = N_0 e^{-\mu x}
\end{equation}
ここで  は比例定数であり、物質の**線吸収係数**（単位: cm-1）と呼ばれる。

また、物質の密度を （単位: g/cm3）とすると、単位面積あたりの質量（面密度） x （単位: g/cm2）を用いて次のように表すこともできる。
\begin{equation}
  N(x) = N_0 e^{-\mu_m \rho_x}
\end{equation}
ここで m=/ は**質量吸収係数**（単位: cm2/g）と呼ばれる。質量吸収係数は物質の密度に依存しないため、線のエネルギーや透過特性を比較・評価する際によく用いられる。

\subsection{放射線の計測}
\begin{figure}[htbp]
  \centering
  \includegraphics[height=10cm, angle=90]{IMG_3846.jpg}
  \caption{GM計数管の構造}
  \label{fig:IMG_3846}
\end{figure}

放射線の検出にはガイガー＝ミュラー計数管（GM計数管）が広く用いられる。GM計数管は、金属筒を陰極、その中心軸に張られた細い金属線を陽極とし、内部にアルゴンなどの不活性ガスと微量の有極性気体（消滅ガス）が封入された構造を持つ（図~\ref{fig:IMG_3846}）。

両極間に高電圧（500〜1000 V 程度）を印加した状態で放射線が管内に入射すると、気体分子が電離されて電子とイオン対が生成される。生じた電子は陽極に向かって加速され、二次電離を繰り返すことで電子雪崩（Townsend discharge）を引き起こす。この結果増幅された脈動電流（パルス）を電気的に検出・カウントすることで、放射線粒子数の計測を行う。なお、入射する放射線のエネルギーの違いによらず一定のパルスが出力されるため、GM計数管はエネルギーの測定ではなく個数の計測に適している。

\begin{figure}[htbp]
  \centering
  \includegraphics[height=10cm, angle=90]{IMG_3847.jpg}
  \caption{放射線計測の実験装置概観}
  \label{fig:IMG_3847}
\end{figure}

\subsection{放射線の単位}
放射線および放射性物質の評価には、その目的に応じて主に以下の単位が用いられる。

\begin{description}
  \item[放射能の単位（ベクレル Bq）]
  単位時間あたりに崩壊する原子核の数を表す単位であり、1秒間に1回の崩壊 を 1Bq=1s-1 と定義する。従来はキュリー（Ci、1Ci=3.71010Bq）が用いられていた。

  \item[吸収線量の単位（グレイ Gy）]
  放射線によって物質 1kg あたりに与えられたエネルギー量を表す物理量であり、1Gy=1J/kg と定義される。

  \item[等価線量・実効線量の単位（シーベルト Sv）]
  吸収線量に放射線の種類（線、線、線など）による生物学的影響の違い（放射線荷重係数）や臓器の感度を考慮した、人体への影響度を表す単位である。
\end{description}

% -----------------------------------------------------------------
\section{装置および測定方法}
\subsection{使用装置}
\begin{itemize}
    \item ガイガー＝ミュラー計数管（GM計数管）およびカウンター
    \item 放射線源（$^{137}\mathrm{Cs}$）
    \item 試料台（GM計数管スタンド）
    \item アルミニウム板（厚さ 1mm）および金属箔 I, II
    \item マイクロメータ
\end{itemize}

\subsection{注意事項}
\begin{itemize}
    \item 実験で使用する放射線源は微量であるが、線源の取り扱いには十分注意する（線源には直接手に触れないこと）。
    \item 線源を紛失することがないよう厳重に注意すること。
    \item GM計数管スタンドの中央（赤色部分）には直接手などが触れないよう注意する。
\end{itemize}

\subsection{測定手順}

\subsubsection{実験1：自然計数の測定}
\begin{enumerate}
    \item 線源を実験台に持ってこない状態で（GM計数管のスタンド内にも入れない）、GM計数管の電源を入れる。
    \item カウンターのゲート時間（計測時間）を 60s に設定する。
    \item 60s 間の計数を 10 回繰り返し、10回の計数値の平均を計算して自然計数の計数値 N0 を求める。
\end{enumerate}

\subsubsection{実験2：放射線の計数値分布の観測}
\begin{enumerate}
    \item 全員が実験1を終えていることを確認する。
    \item 線源を受け取り（記録簿に必要事項を記入）、GM計数管スタンドの下段に入れる。
    \item カウンターのゲート時間を 1s に設定し、計数値を記録する（最初に100回、次に300回繰り返す）。
    \item 線源を入れる位置を変えて、300回の計測を行う。
\end{enumerate}

\subsubsection{実験3：線の吸収係数の観測}
\begin{enumerate}
    \item GM計数管スタンドに線源を置き、線源とGMカウンターの間の穴の開いた板を入れる。線源の位置を記録し、この後の実験では線源の位置を変えないこと。
    \item 穴の開いた板の上に、厚さ 1mm のAl板を置き、1s 間の計数値 N' を記録する。計測は複数回行い、平均値 N' を求めておく。
    \item 1mm のAl板を取り出し、金属Iの薄板を0枚から1枚ずつ入れて、各枚数に対して計数値 N を複数回測定し、平均値 N を求める。
    \item 金属Iの薄板の厚さをマイクロメータで計測する。1枚の薄板に対して測定部分を変えて複数の計測を行い、その平均を求める。さらに複数枚の厚みに対して同様に平均値を求め、それらを平均して1枚の厚さの平均値を求める（マイクロメータのゼロ点が入っている場合はゼロ点補正を行うこと）。
    \item 金属IIについて同様の測定を行う。
\end{enumerate}

% -----------------------------------------------------------------
\section{実験結果}
測定時の環境条件は以下の通りである。
\begin{itemize}
    \item 測定日：2026年10月17日 % ※実際の年に合わせてください
    \item 気温：$24\,^\circ\mathrm{C}$
    \item 気圧：1019hPa
\end{itemize}

\subsection{実験1：自然計数の測定}

測定を10回行った際の各回の計数値、およびその平均値（自然計数の計数値 N0）を表~\ref{tab:natural_count_raw}に示す。また、担当教員の指示により測定回数は10回とし、その計数値 N0 は 16.1cpm となった。

\begin{table}[htbp]
    \centering
    \caption{自然計数の測定結果（生データ）}
    \label{tab:natural_count_raw}
    \begin{tabular}{cc}
        \hline
        回数 / 回 & 計数値 N / cpm \\
        \hline
        1 & 16 \\
        2 & 17 \\
        3 & 14 \\
        4 & 16 \\
        5 & 15 \\
        6 & 19 \\
        7 & 17 \\
        8 & 14 \\
        9 & 22 \\
        10 & 11 \\
        \hline
        平均 (N0) & 16.1 \\
        \hline
    \end{tabular}
\end{table}
\clearpage
さらに、レポートの課題に基づき、同じ計数値 N ごとの出現回数 n をまとめ直したヒストグラムを図~\ref{fig:natural_count_hist}に示す。

\begin{figure}[htbp]
    \centering
    \includegraphics[width=0.6\linewidth]{hist.png}
    \caption{自然計数の出現頻度}
    \label{fig:natural_count_hist}
\end{figure}
\subsection{実験2：放射線の計数値分布の観測}

担当教員の指示により、本実験では線源の位置を変える測定や300回の測定は行わず、ゲート時間 1s における総測定回数100回のみのデータ取得を行った。

得られた100回のデータについて各計数値の出現回数と確率、および平均値 N=8.32 を用いたポアソン分布の理論値 P(N) などを整理した結果を表~\ref{tab:exp2_data}に示す。

\begin{table}[htbp]
    \centering
    \caption{計数値の分布の実験データ（測定回数100回,ゲート時間1秒,線源の位置60mm）}
    \label{tab:exp2_data}
    \resizebox{\textwidth}{!}{
    \begin{tabular}{cccccc}
        \hline
        計数値 & 出現回数 & 出現確率 & 計数値とnNの積 & 2乗偏差とnNの積 & ポアソン分布 \\
        N & nN & nN/nN & nNN & nN(N-N)2 & P(N) \\
        \hline
        1 & 1 & 0.01 & 1 & 53.58 & 0.002 \\
        3 & 2 & 0.02 & 6 & 56.60 & 0.023 \\
        4 & 7 & 0.07 & 28 & 130.64 & 0.049 \\
        5 & 9 & 0.09 & 45 & 99.20 & 0.081 \\
        6 & 7 & 0.07 & 42 & 37.68 & 0.112 \\
        7 & 13 & 0.13 & 91 & 22.65 & 0.133 \\
        8 & 16 & 0.16 & 128 & 1.64 & 0.139 \\
        9 & 12 & 0.12 & 108 & 5.55 & 0.128 \\
        10 & 9 & 0.09 & 90 & 25.40 & 0.107 \\
        11 & 11 & 0.11 & 121 & 79.01 & 0.081 \\
        12 & 7 & 0.07 & 84 & 94.80 & 0.056 \\
        13 & 1 & 0.01 & 13 & 21.90 & 0.036 \\
        14 & 2 & 0.02 & 28 & 64.52 & 0.021 \\
        15 & 1 & 0.01 & 15 & 44.62 & 0.012 \\
        16 & 2 & 0.02 & 32 & 117.96 & 0.006 \\
        \hline
        合計  & 100 & 1.00 & 832 & 855.76 & - \\
        \hline
    \end{tabular}
    }
\end{table}

表~\ref{tab:exp2_data}より、本データの平均値 N および標準偏差  は以下の通り求められた。
\begin{align*}
    \bar{N} &= \frac{\sum n_N N}{\sum n_N} = \frac{832}{100} = 8.32 \\
    \sigma &= \sqrt{\frac{\sum n_N (N - \bar{N})^2}{\sum n_N - 1}} = \sqrt{\frac{855.76}{99}} \approx 2.94
\end{align*}
また、ポアソン分布において標準偏差の期待値となる N を計算すると、8.322.88 であった。N の関係がよく成り立っていることがわかる。

さらに、表~\ref{tab:exp2_data}のデータを用いて作成した、計数値の出現確率のヒストグラムおよびポアソン分布の理論値 P(N) （黒丸プロット）の比較を図~\ref{fig:poisson_dist}に示す。

\begin{figure}[htbp]
    \centering
    \includegraphics[width=0.7\linewidth]{hist2.png} 
    \caption{計数値の出現確率とポアソン分布の比較}
    \label{fig:poisson_dist}
\end{figure}
% -----------------------------------------------------------------
\newpage
\subsection{β線の吸収の測定}
金属I（Cu箔）および金属II（Al箔）について、透過する線の計数値を測定した結果をそれぞれ表~\ref{tab:exp3_data_cu}および表~\ref{tab:exp3_data_al}に示す。

各計数値については、測定された計数値の平均値 N から自然計数（バックグラウンド）の平均値 N'=16.1cpm を差し引くことで、正味の線計数値 N を求めた。
また、表中の金属箔の厚さ x は、マイクロメータを用いて1枚の箔につき5箇所の厚みを計測してその平均値を求め、そこに重ねた枚数を掛けることで算出したものである。
\begin{table}[htbp]
    \centering
    \caption{線の吸収の実験データ（金属I：Cu箔）}
    \label{tab:exp3_data_cu}
    \resizebox{\textwidth}{!}{
    \begin{tabular}{ccccccccc}
        \hline
        枚数 & 厚さ x & 1分間の計数値 & 平均値 & 線計数値 & 対数 & 標準偏差 & \multicolumn{2}{c}{標準偏差の範囲} \\
         & /mm & N & N & N & eN &  & e(N-) & e(N+) \\
        \hline
        0 & 0.000 & 502, 524, 485, 495, 487 & 498.6 & 482.5 & 6.179 & 10.1 & 6.158 & 6.200 \\
        1 & 0.236 & 316, 329, 285, 304, 291 & 305.0 & 288.9 & 5.666 & 8.0 & 5.638 & 5.693 \\
        3 & 0.708 & 203, 214, 209, 213, 221 & 212.0 & 195.9 & 5.278 & 6.8 & 5.243 & 5.311 \\
        5 & 1.180 & 138, 140, 164, 146, 124 & 142.4 & 126.3 & 4.839 & 5.6 & 4.793 & 4.882 \\
        \hline
    \end{tabular}
    }
\end{table}

\begin{table}[htbp]
    \centering
    \caption{線の吸収の実験データ（金属II：Al箔）}
    \label{tab:exp3_data_al}
    \resizebox{\textwidth}{!}{
    \begin{tabular}{ccccccccc}
        \hline
        枚数 & 厚さ x & 1分間の計数値 & 平均値 & 線計数値 & 対数 & 標準偏差 & \multicolumn{2}{c}{標準偏差の範囲} \\
         & /mm & N & N & N & eN &  & e(N-) & e(N+) \\
        \hline
        0 & 0.000 & 502, 524, 485, 495, 487 & 498.6 & 482.5 & 6.179 & 10.1 & 6.158 & 6.200 \\
        1 & 0.282 & 463, 439, 439, 430, 466 & 447.4 & 431.3 & 6.067 & 9.6 & 6.044 & 6.089 \\
        3 & 0.846 & 378, 314, 334, 374, 375 & 355.0 & 338.9 & 5.826 & 8.6 & 5.799 & 5.851 \\
        5 & 1.410 & 296, 307, 293, 268, 293 & 291.4 & 275.3 & 5.618 & 7.8 & 5.589 & 5.646 \\
        \hline
    \end{tabular}
    }
\end{table}
\section{考察}
ここに考察を記述します。

% -----------------------------------------------------------------
\section{結論}
ここに結論を記述します。

% -----------------------------------------------------------------
\begin{thebibliography}{9}
  \bibitem{text} 著者名, 『文献名』, 出版社
\end{thebibliography}
\end{document}
